% !TeX root = main.tex
\newpage
\section{Photon Lighting Analysis}

\subsection{Current LightingBall State}

The current \texttt{LightingBall} implementation has crossed from
photon-tracer bring-up into the first surface-lighting system.  Photon
particles are generated as mobile runtime particles with photon \texttt{ptype}.
They move through the same source-owned compute schedule as other mobile
particles, but their optical rule is special:
\begin{itemize}
    \item photons ignore other photons;
    \item regular particles ignore photon targets;
    \item photon-source contacts with regular particles, walls, or the
          lighting sphere are optical events rather than mechanical contact
          resolution;
    \item photon speed is preserved by reflecting the velocity direction;
    \item \texttt{colFlg} marks the photon as in collision for the render
          frame;
    \item photon \texttt{material\_id} is temporarily changed to the contacted
          surface material for diagnostics and material-source checks.
\end{itemize}

This is useful for debugging whether photons hit the correct surface and
whether the material color source is being read correctly.  Surface lighting
itself is now owned by \texttt{BoundaryLightRecord} state in the
\texttt{ParticleLighting} Vulkan path.  For \texttt{LightingBall}, photon
contact deposits the contacted material color into the spatial cell containing
the analytic hit point, and the lighting surface fragment shader reads that
cell-owned color.

\subsection{Why Direct Photon Rendering Was Unsatisfying}

Direct photon rendering makes sparse sampling visible.  A small number of
photons produces isolated points, streaks, or momentary contact flashes.  Even
with denser emission, the render result is tied to where the photon particles
happen to be at the current frame rather than to accumulated light on the
surface.

The current \texttt{LUMENS} path also mixes two concepts that should be
separate:
\begin{itemize}
    \item photon tracer visibility: whether and how a photon particle itself
          should be drawn for debugging;
    \item surface illumination: how much light a boundary receives and later
          displays.
\end{itemize}

The eye-alignment brightness calculation belongs to the optional tracer path.
It is not a surface-lighting equation because moving the camera or eye should
not change how much light was deposited on a wall or sphere.

\subsection{Target Direction From PhotonLighting.tex}

\texttt{PhotonLighting.tex} describes the right target architecture:
photons are transport samples, while renderable boundaries accumulate and
display received light.  The intended pipeline is:
\begin{enumerate}
    \item move photon particles at fixed speed;
    \item detect photon-boundary intersections;
    \item compute hit position and boundary normal;
    \item compute incidence-weighted deposited energy;
    \item accumulate that energy into a boundary-light buffer;
    \item optionally spread the deposit over nearby boundary samples;
    \item temporally filter the boundary-light buffer;
    \item render the boundary from sampled boundary-proxy color;
    \item optionally render a sampled subset of photons as tracers.
\end{enumerate}

The important shift is that photons no longer need to be dense enough to fill
the visible surface.  They only need to sample light transport well enough for
the boundary-light buffer to converge to a usable illuminated image.

\subsection{What We Already Have}

Several pieces of the target system already exist in partial form:
\begin{itemize}
    \item photon generation through \texttt{particle\_streams};
    \item stable photon behavior through runtime \texttt{ptype};
    \item fixed-speed photon reflection;
    \item photon-wall, photon-sphere, and photon-particle optical contact
          detection;
    \item surface material color lookup through \texttt{material\_id};
    \item analytic sphere normals in the lighting-ball evaluator;
    \item rectangle-wall normals from \texttt{rectangle\_wall\_segments};
    \item material debug visibility and debug color for optional photon
          tracers.
\end{itemize}

These pieces are enough to prove transport, contact, and material-source
behavior.  The current Vulkan \texttt{LightingBall} path adds the first
persistent illumination owner by storing color in one lighting record per
spatial cell.

\subsection{Remaining Missing Pieces}

The missing architecture is no longer the basic boundary-owned light state.
That first color path is now present in Vulkan as cell-owned
\texttt{BoundaryLightRecord} state.  The remaining gap is richer light
response, especially spectral lighting.

The next target should answer what additional state is required beyond
\texttt{rgb\_valid}.  Candidate additions include:
\begin{itemize}
    \item photon wavelength, band, or spectral bin;
    \item material spectral response or absorption;
    \item separate debug views for deposited color versus spectral state;
    \item optional accumulated direction/coherence for later specular work;
    \item debug counters for active photons, hits, and cell writes.
\end{itemize}
Atomics, additive current-frame accumulation, and temporal filtering are not
part of the first cell-owned color implementation.  They remain possible later
extensions after spectral requirements are defined.

\subsection{Recommended Migration Path}

The safest path is incremental and keeps the current photon tracer behavior as
a diagnostic while surface lighting is built.

\subsubsection{Stage 1: Cell-Owned Boundary Color}

Use the existing photon-boundary contact events.  For each photon hit, compute:
\[
    C_{\mathrm{incidence}} =
    \max(-\mathbf d_p \cdot \mathbf n_b, 0)
\]
and deposit:
\[
    E_{\mathrm{deposit}} =
    E_p C_{\mathrm{incidence}} .
\]

The active Vulkan implementation writes exactly one spatial cell.  For the
analytic sphere, the contact normal gives the hit point:
\begin{verbatim}
LIGHTING_BALL_CENTER + contactNormal * LIGHTING_BALL_RADIUS
\end{verbatim}
The cell containing that point stores the contacted material color until
another photon overwrites it.  This verifies the sign of the normal, the
material color, and whether writes occur on the expected object.

\subsubsection{Stage 2: Render Boundary-Light Values}

Add a boundary render path that reads accumulated light and applies it to the
surface base color:
\[
    \mathbf C_{\mathrm{final}} =
    \mathbf C_{\mathrm{base}}
    \left(A_{\mathrm{ambient}} + G\mathbf L_{\mathrm{filtered}}\right).
\]

At this stage the result may be speckled.  That is acceptable.  The goal is to
prove that the boundary, not the photon point, is emitting the visible light.

\subsubsection{Stage 3: Spatial Boundary-Space Spread}

This earlier proxy-kernel stage is superseded for the current Vulkan
\texttt{LightingBall} path.  The active renderer colors the analytic sphere
from the owning cell id carried by the render vertex and read in the fragment
shader.  A kernel may still be useful later for filtering or smoothing, but it
is not the first visible surface-lighting rule.

The retired proxy-kernel form was:
\[
    w_j =
    \begin{cases}
        (1-r_j/R_L)^2, & r_j < R_L,\\
        0, & r_j \ge R_L.
    \end{cases}
\]

Normalize the weights so increasing the number of nearby proxies does not
increase brightness.

\subsubsection{Stage 4: Optional Temporal Or Additive Light}

Only after persistent cell-owned color is stable should the model consider
current-frame additive accumulation, temporal filtering, or decay.  Those
features require extra state and, in Vulkan, may require atomics or ping-pong
buffers.  They are deliberately outside the first boundary-space
implementation.

\subsubsection{Stage 5: Material Optical Response}

Add optical material controls after deposition is visible:
\begin{itemize}
    \item absorption;
    \item reflection;
    \item roughness;
    \item photon energy reduction or termination.
\end{itemize}

Until this stage, boundaries can be treated as fully absorbing for deposition
while still reflecting photons for transport debugging.

\subsubsection{Deferred Directional Specular}

Photon-direction specular lighting is a useful later extension.  A boundary
sample could accumulate the incoming photon direction in addition to deposited
energy, and a coherent set of incoming directions could produce a highlight
from the photon simulation itself.  This is intentionally deferred until the
first scalar path is proven:
\[
    \textrm{photon hit}
    \rightarrow
    \textrm{boundary sample energy}
    \rightarrow
    \textrm{filtered lit sector}.
\]

\subsection{Boundary Sample Choice}

For \texttt{LightingBall}, the first Vulkan lighting owner is the spatial cell
that contains the analytic sphere hit point.  Existing generated sphere markers
remain useful because they make nearby analytic sphere geometry discoverable,
but they no longer own the final Vulkan surface color.

A separate spectral or energy sample set is deferred.  The current direction is
to learn from the cell-owned color path first, then decide what extra Vulkan
tables are required for spectral lighting.

\subsection{Effect On Current Debug Controls}

\texttt{debug\_visible} and \texttt{debug\_color} should remain photon-tracer
controls.  They answer whether uncontacted photons should be drawn so their
travel can be inspected.  They should not affect deposited light.

\texttt{gl\_point\_size} also remains a tracer/render diagnostic.  It changes
the apparent size of photon points but should not change the boundary-light
deposition radius.  The deposition radius belongs to a separate lighting
configuration value such as \texttt{photon\_light\_influence\_radius}.

\subsection{Near-Term Decision}

The current decision is to use persistent cell-owned boundary-space color for
Vulkan \texttt{LightingBall}.  The cleaned scene removes the old boundary-light
sample/state path and visible surface-cell path, and the current Vulkan
implementation proves:
\[
    \textrm{photon hit}
    \rightarrow
    \textrm{persistent cell-owned surface color}
    \rightarrow
    \textrm{analytic surface fragment color}.
\]

The next decision is how spectral lighting should extend this cell-owned
resource.
